\section{带标签的提升：计算真正保留下来的母方向}

设 \(C\subset S^{D-1}\) 是一个接吻构型，将目标空间分解为
\(\mathbb R^D\oplus\mathbb R^k\)。选择两两不同、非零且范数小于一的尾向量
\(t_i\)，并为每个尾标签选择母方向集合 \(A_i\subseteq C\)。记

\[
a_i=\sqrt{1-\|t_i\|^2},\qquad U=\bigcup_iA_i.
\]

保留 \(x\in C\setminus U\) 对应的赤道点 \((x,0)\)，将 \(x\in A_i\)
提升为 \((a_ix,t_i)\)，再加入纯尾接吻码 \(Y\subset S^{k-1}\)。总点数为

\[
|X|=|C|-|U|+\sum_i|A_i|+|Y|.
\]

同一个母方向无论被多少标签复用，在赤道上都只删除一次。对不同的带标签点
\((x,i)\ne(x',j)\) 及纯尾点 \(y\in Y\)，还需验证

\[
a_ia_j\langle x,x'\rangle+\langle t_i,t_j\rangle\le\frac12,
\qquad
\langle t_i,y\rangle\le\frac12.
\]

\subsection{复用次数与四三角形构造}

如果全部尾向量的平方范数均为 \(h^2<1/2\)，同一母方向的不同副本必须满足

\[
\langle t_i,t_j\rangle\le h^2-\frac12.
\]

若出现 \(m\) 个副本，其尾向量之和的平方范数至多为

\[
\frac m2\bigl(1-m(1-2h^2)\bigr).
\]

非负性给出纬度界。尾向量的互内积严格为负，还迫使它们仿射无关：
否则两个不相交子集的正线性组合相等，但二者的内积却为负。因此
\[
m\le r=\min\{k+1,\lfloor(1-2h^2)^{-1}\rfloor\}.
\]
正单纯形尾向量使单个方向达到此界。若第 $i$ 个标签容量为 $\alpha_i$，
则总点数至多为 $|C|+|Y|+(1-1/r)\sum_i\alpha_i$。

更精确地，将非空块 $B$ 的每个方向都放在同一组 $r\ge2$ 个尾标签上，
并允许选择 $0<h^2<1/2$，其充要条件是 $k\ge r-1$ 且块内不同方向内积
至多为 $1/(r+1)$。必要性来自复用次数界及随 $h^2$ 递减的同标签阈值
$(1/2-h^2)/(1-h^2)$；取 $h^2=(r-1)/(2r)$ 的正单纯形即得充分性，
无需纯尾点便有 $|C|+(r-1)|B|$ 点。当 $h^2=1/3$ 时，等号对应零和正三角形。

将 \(D_3\) 的十二个根缩放至尾范数
\(1/\sqrt3\)，可以把它们划分成四个零和三角形。四个互不相交的 496 点
Leech 块分别使用一个三角形，得到

\[
196560-4(496)+12(496)+12=200540.
\]

原来的分配复用次数为 \((3,3,2,2,2)\)，新分配为
\((3,3,3,3)\)。提升层点数没有改变，但少删除了 496 个赤道点。496
点块来自公开的 PackingStar 数据；这里改变的是它们与尾标签之间的分配。

对这个固定的十二标签模型，若每个标签最多容纳 \(\alpha\)
个母方向，则总点数至多为 \(196560+8\alpha+12\)。四三角形构造在
\(\alpha=496\) 时达到这一容量界。

\subsection{后续的 201,567
点构造}

后续构造扩展公开的双层提升方案。第一层沿用 Lindow 的 201,566
点贡献中的数据：2,258 个平方范数为 192 的整数头部，其中 2,090
个各删除一个 Leech 最短向量，168
个不需要删除赤道点。每个头部产生三个帽点。

本节采用 Leech 最短向量平方范数为 32、接吻内积阈值为 16
的整数约定。第一层头部 \(y\) 产生

\[
(y/3,\;4v/\sqrt3),
\]

其中 \(v\) 遍历所分配的 \(D_3\)
三角形。相同三角形上的两个不同头部的整数内积至多为
48，不同三角形之间至多为 96。

第二层选择 Leech 最短向量 \(u\)
作为基点，并将其分配到三个坐标直线之一。删除赤道点 \((u,0)\)，加入

\[
(\sqrt3\,u/2,\;\pm\sqrt8\,e_j).
\]

新的见证共有 311 个基点，三个标签分别有 105、103、103
个。它们互不重复，均不在第一层的删除集合中；与每个第一层头部的内积满足
\(y\cdot u\le32\)；同一标签上的两个基点内积至多为八，不同标签之间至多为十六。

例如，第一层与第二层的内积上界为

\[
\frac{\sqrt3}{6}\,32+\frac{4\sqrt8}{\sqrt3}
=\frac{16\sqrt3+8\sqrt6}{3}<16.
\]

第二层与纯尾点之间的上界为
\(\sqrt8\sqrt{32}=16\)，可以达到等号。同一基点的两个帽点内积也为十六。逐类验证后，点数为

\[
196560-2090-311+3(2258)+2(311)+12=201567.
\]

将公开输入的 310 个第二层基点换成新选出的 311
个基点，便净增一个点：每个基点用两个帽点替换一个赤道点。
